\setcounter{chapter}{1}

\chapter{Related Work}\label{sec:relatedwork}


Controllable narrative generation is increasingly central to interactive media. Our work builds on an existing interactive narrative system developed at Disney Research Studios, extending it to handle real-time interruptions and time pressure. To understand this challenge, we examine how interactive narrative systems have evolved, from symbolic planning approaches to modern LLM-based generation. Our focus narrows to two key areas: multi-agent architectures and beat-based representations that are most relevant to controllable storytelling under interruption and time constraints.

\section{Symbolic and Goal-Driven Narrative Planning}

Early narrative systems relied on symbolic planning, treating story generation as a search problem through logical states and actions \cite{Meehan1977}. Systems like IPOCL \cite{riedl2010narrative} and Sabre \cite{ware2021sabre} extended this by modelling character intentions and beliefs, so that plans stay causally coherent while characters act believably. These approaches guaranteed logical coherence by construction, but they produced rigid, constrained narratives that lacked the fluidity interactive applications demanded.

A foundational insight emerged from {Façade} \cite{mateas2003facade}, a landmark interactive drama system. Its key lesson: narrative control becomes more robust when you separate what should happen next (narrative deliberation) from how characters express it (performance and dialogue generation). The system's three-layer architecture operated with story planning on one level, character behavior on another, and performance at a third. This decoupling enabled interactive steering while preserving character believability. The principle remains influential today: separating narrative decisions from character-level execution lets systems respond to user input without losing coherence. Every hierarchical approach that has followed builds on this insight.

\section{LLM-based Storytelling}\label{sec:rw-llm-story}

The emergence of large language models revolutionized narrative generation by replacing symbolic planning with learned, fluent text generation. However, unguided LLMs suffer from narrative drift \cite{yang2023doc}. They meander without maintaining explicit plot structure, contradict earlier plot points, or fail to reach satisfying conclusions. This limitation motivated a return to hierarchical architectures that combine LLM generative power with explicit narrative control.


Recent work has explored multi-agent frameworks in which specialized LLM agents coordinate narrative generation.
CoDi \cite{kim2025codi} exemplifies this approach. It uses a multi-agent framework with four specialized agents that work together:

\begin{itemize}
    \item \textbf{Planner LLM}: Sets up the story world, character profiles, and high-level narrative goals.
    \item \textbf{Director LLM}: Reads the story so far each turn and issues a high-level instruction of one of four types (character action, description, intervention, or ending) to advance those goals.
    The narrative goals are encoded as a utility function: a list of natural-language goals, each with a numerical weight reflecting its importance. The story state is not tracked explicitly; the director infers it from the narrative so far.
    \item \textbf{Character LLM} (one per character): Generates the character's inner thought, action and dialogue in response to the director's instructions, guided by its own profile and utility function.
    \item \textbf{Editor LLM}: Formats the output into screenplay or novel form.
\end{itemize}

CoDi shows that giving the director explicit narrative goals and dedicated intervention instructions yields stories that outperform a prior director-actor system and approach human-written stories on several dimensions. Yet CoDi is designed for adaptive storytelling, where the system flexibly adjusts narratives to achieve predefined goals. Its evaluation uses autonomously generated stories without live user input, and supporting more complex user interactions is left to future work. This is a critical distinction. CoDi's character agents periodically update their own utility functions as the story unfolds, but the paper describes no mechanism for detecting a user interruption or steering back to narrative coherence.

\section{Beat-Based Narrative Representation}

While CoDi represents narrative goals as weighted utility functions, an alternative approach represents stories as sequences of discrete narrative beats. Each beat is a distinct moment or scene with clear dramatic function. This abstraction is grounded in screenwriting theory \cite{mckee1997story} and provides a middle ground between rigid symbolic planning and unconstrained LLM generation.

NarrativeGenie \cite{kumaran2024narrativegenie} is the most directly relevant precedent in this space. It converts free-text story outlines into structured sequences of narrative beats, establishes causal dependencies between beats to form a partial order, and grounds characters and props in the world. Dialogue and actions are then generated within each beat. By constraining generation to a single beat at a time, the system maintains coherence throughout the story. It also offers a degree of interruption handling: when a player acts out of order, it generates alternate character dialogue intended to guide the player back on track, and an adaptive dialogue manager responds to the current game state in real time.

Beats offer three key advantages over utility-function-based approaches. First, interpretability: beat objectives are expressed in natural language, making narrative logic transparent and debuggable, unlike the opaque numerical trade-offs of utility functions. Second, checkpointing: each beat represents a natural checkpoint where the system can verify whether the goal is satisfied, enabling principled decisions about recovery and pacing. Third, temporal reasoning: knowing N beats remain and t seconds remain means we can compute a pacing target directly, providing a clearer implementation path than utility functions allow.

This thesis builds on beats as its core story representation and relies in particular on checkpointing and temporal reasoning: beats give the Director a clear point to steer back to after an interruption, and a count of remaining beats to pace the story against a time limit. In Section~\ref{sec:llm-beats-ablation}, we test whether this representation helps by comparing it with a Director that is given the full story text instead.

\section{Interruption Handling }

Several interactive narrative systems demonstrate approaches to handling unexpected user inputs. ``Bounded autonomy'' \cite{Guo2026} is a control approach for LLM-driven characters in a live multiplayer social game: a ``whisper'' channel lets players softly steer a character's next action without overriding its autonomy, while action grounding keeps behaviour within what the game can execute. AI Dungeon \cite{hua2020playing}, a text adventure game built on a language model, generates story responses to free-form player input in real time, allowing players to interrupt and redirect the narrative dynamically. labs.character.ai's experimental books format \cite{characterai2026labs} lets users play through public-domain books as a character, follow the story, or rewrite it. Its handling of interruptions is not publicly documented. In an informal test, we played Holmes in a scene from \textit{The Adventure of the Speckled Band} and answered every character turn with the same unrelated statement (``I'm going to Italy''), 14 times in a row. The characters acknowledged each interruption briefly, for example by telling Holmes that the trip could wait, and returned to the scene, so the story did not derail. However, the plot also barely advanced: Helen Stoner restated her plea for help in new words turn after turn but never told her story, and the scene only moved on when a new event (a caller at the door) was introduced. In short, the user is acknowledged, but the plot does not move forward.

These systems maintain narrative control through context windows and character profiles that constrain generation. However, the challenge is not merely responding to interruptions, but doing so while maintaining long-term plot coherence and meeting temporal deadlines.

\section{Time-Constrained Generation}

A distinct but crucial challenge is narrative completion within fixed time budgets. Prior work on time-aware planning \cite{porteous2011controlling,CavazzaCharlesMead2001} has addressed temporal constraints in symbolic storytelling and embodied agents, but systematic treatment of pacing adaptation in LLM-based systems remains largely unexplored. Most LLM systems ignore time entirely, relying on implicit and unreliable length targets. Our work addresses this gap by introducing explicit temporal monitoring mechanisms.

\section{LLM-as-Judge}

Assessing narrative quality is inherently subjective: it spans coherence, character consistency, recovery from interruption and appropriate pacing.

Recent work has therefore adopted LLM-as-Judge evaluation \cite{zheng2023judging}, which CoDi \cite{kim2025codi} also uses. Structured prompts let a language model score transcripts on several dimensions without extensive manual annotation. This scales to many samples and makes it possible to score complex, subjective properties such as how well a system recovers from user interruptions.

LLM judges are not neutral, however. In a systematic study of this approach (MT-Bench and Chatbot Arena), most judges were biased towards the answer shown first, and only GPT-4 gave consistent verdicts in more than 60\% of cases \cite{zheng2023judging}.

We therefore judge every pair of transcripts in both orders (AB and BA) and only count a winner when both orders agree. Because an LLM judge may still miss what users actually experience, we complement it with a user study. 

